\chapter[1960 --- Burnet et al.]{1960: Frank Macfarlane Burnet, Peter Brian Medawar}
\label{ch:1960_Burnet}

\section*{Main Contribution}

\textit{Discovered acquired immunological tolerance---the ability of the immune system to be trained to accept foreign tissues---the basis for organ transplantation and autoimmune disease research.}

\section{Official Citation}

for their discovery of acquired immunological tolerance

\section{Background and Historical Context}

The 1960 Nobel Prize in Physiology or Medicine was awarded to Frank Macfarlane Burnet and Peter Brian Medawar for their discovery of acquired immunological tolerance. This discovery fundamentally changed our understanding of the immune system and laid the foundation for modern transplantation medicine. The work of Burnet and Medawar demonstrated that the immune system could be educated to tolerate foreign tissues, a concept that was essential for the development of organ transplantation and the treatment of autoimmune diseases.

\subsection{The Immunology Revolution}

By the mid-20th century, immunology had emerged as a major field of biological research. Researchers had identified the major components of the immune system, including antibodies, white blood cells, and the lymph nodes and spleen. However, the mechanisms by which the immune system distinguished between self (the body's own tissues) and non-self (foreign tissues and pathogens) were still poorly understood.

The discovery of acquired immunological tolerance by Burnet and Medawar provided a new framework for understanding how the immune system makes this distinction. Their work demonstrated that the immune system could be "educated" during early life to tolerate foreign tissues, opening up new possibilities for the treatment of transplant rejection and autoimmune diseases.

\subsection{Frank Macfarlane Burnet}

Frank Macfarlane Burnet was born on September 3, 1899, in Traralgon, Australia. He studied medicine at the University of Melbourne, where he developed an interest in virology and immunology. After completing his medical degree, Burnet worked at the Walter and Eliza Hall Institute of Medical Research in Melbourne, where he would spend most of his career.

Burnet was a innovative researcher who made numerous contributions to immunology and virology. He is best known for his clonal selection theory, which provided a comprehensive explanation for how the immune system responds to foreign antigens.

\subsection{Peter Brian Medawar}

Peter Brian Medawar was born on February 28, 1915, in Petropolis, Brazil. He studied at the University of Oxford, where he earned his doctorate in 1937 for his work on the biology of wound healing. After World War II, Medawar joined the faculty at the University of Birmingham and later at University College London, where he would spend most of his career.

Medawar was a charismatic and persuasive researcher who was skilled at designing experiments and interpreting results. His work on acquired immunological tolerance would establish the foundations of modern transplantation medicine.

\subsection{The Problem of Transplant Rejection}

By the mid-20th century, it had become clear that the immune system played a major role in the rejection of transplanted tissues. When tissues from one individual were transplanted into another, the recipient's immune system would recognize the transplanted tissue as foreign and mount an immune response against it, leading to the destruction of the transplanted tissue.

This problem was a major obstacle to the development of organ transplantation. Transplant surgeons could perform the technical aspects of transplantation, but they had no way to prevent the immune system from rejecting the transplanted organ. The discovery of acquired immunological tolerance by Burnet and Medawar would provide a solution to this problem.

\section{The Discovery/Contribution}

The work of Burnet and Medawar on acquired immunological tolerance involved the development of a new theory of immunity and the experimental demonstration that the immune system could be educated to tolerate foreign tissues.

\subsection{Burnet's Clonal Selection Theory}

Frank Macfarlane Burnet's primary contribution to immunology was the development of the clonal selection theory, which provided a comprehensive explanation for how the immune system responds to foreign antigens.

\subsubsection{The Theory}

The clonal selection theory, which Burnet first proposed in 1957, states that:

\begin{itemize}
\item The immune system contains a vast repertoire of lymphocytes, each with a unique receptor that can recognize a specific antigen
\item When a lymphocyte encounters an antigen that matches its receptor, it becomes activated and proliferates, producing a clone of identical cells
\item These activated lymphocytes (or their descendants) produce antibodies or carry out cell-mediated immune responses against the antigen
\item Lymphocytes that recognize self-antigens are eliminated or inactivated during early development, preventing autoimmune responses
\end{itemize}

\subsubsection{Significance of the Theory}

The clonal selection theory was significant because it explained several key features of the immune response:

\begin{itemize}
\item It explained the specificity of the immune response---why different antigens stimulate the production of different antibodies
\item It explained the memory of the immune response---why the response to a previously encountered antigen is faster and stronger than the initial response
\item It explained tolerance---how the immune system can be educated to tolerate foreign tissues during early development
\end{itemize}

\subsection{Medawar's Experimental Demonstration of Tolerance}

Peter Brian Medawar's primary contribution was the experimental demonstration that acquired immunological tolerance could be induced in animals. Medawar's experiments provided direct evidence for Burnet's theory and laid the foundation for modern transplantation medicine.

\subsubsection{The Experimental Design}

Medawar's experiments involved injecting cells from one strain of mice into newborn mice of a different strain. When these mice grew up, they were unable to mount an immune response against skin grafts from the donor strain, demonstrating that they had become tolerant to the donor tissue.

\subsubsection{Key Results}

Medawar's key results included:

\begin{itemize}
\item Injection of donor cells into newborn mice could induce tolerance to donor tissue
\item The tolerance was specific---the mice were tolerant to donor tissue but could still mount immune responses against tissue from other strains
\item The tolerance was permanent---once established, it persisted for the lifetime of the animal
\end{itemize}

\subsubsection{Mechanism of Tolerance}

Burnet and Medawar proposed that tolerance was established during early development by the elimination or inactivation of lymphocytes that recognized foreign antigens. This process, which they called clonal deletion, was thought to occur during a critical period of immune system development when the immune system was being "educated" to distinguish self from non-self.

\subsection{The Combined Contribution}

The work of Burnet and Medawar was complementary and synergistic. Burnet provided the theoretical framework (the clonal selection theory) that explained how tolerance could occur, while Medawar provided the experimental evidence that demonstrated tolerance in living animals. Together, their work established the concept of acquired immunological tolerance and laid the foundation for modern transplantation medicine.

\section{Impact on Modern Medicine}

The discovery of acquired immunological tolerance by Burnet and Medawar has had a significant impact on modern medicine. Their work has contributed to the development of organ transplantation, the treatment of autoimmune diseases, and our understanding of the immune system.

\subsection{Organ Transplantation}

The most immediate impact of the work of Burnet and Medawar was in the field of organ transplantation. The discovery that tolerance could be induced in animals suggested that it might be possible to induce tolerance in human patients receiving organ transplants. This idea led to the development of immunosuppressive drugs and other strategies for preventing transplant rejection.

Today, organ transplantation is a routine procedure that saves thousands of lives each year. The immunosuppressive drugs used to prevent transplant rejection are based on the principles established by Burnet and Medawar, and the development of new transplantation strategies continues to build on their work.

\subsection{Autoimmune Diseases}

The clonal selection theory and the concept of tolerance have also contributed to our understanding of autoimmune diseases. Autoimmune diseases occur when the immune system fails to recognize self-antigens and mounts an immune response against the body's own tissues. The work of Burnet and Medawar has provided insights into the mechanisms of autoimmune disease and has led to the development of new treatments.

\subsection{Immunodeficiency Diseases}

The clonal selection theory has also contributed to our understanding of immunodeficiency diseases---conditions in which the immune system is unable to mount effective immune responses against pathogens. The theory has helped researchers understand the mechanisms of immunodeficiency and has led to the development of new treatments.

\subsection{Vaccine Development}

The clonal selection theory has also contributed to our understanding of how vaccines work. Vaccines work by stimulating the immune system to produce a memory response against specific pathogens, so that the immune system can respond quickly and effectively if the pathogen is encountered again. The principles established by Burnet and Medawar have been essential for the development of modern vaccines.

\subsection{Tolerance-Based Therapies}

More recently, the concept of acquired immunological tolerance has inspired the development of new therapies for autoimmune diseases and transplant rejection. These include:

\begin{itemize}
\item Tolerogenic vaccines: Vaccines designed to induce tolerance to specific antigens, rather than stimulating an immune response against them
\item Regulatory T-cell therapy: The use of regulatory T cells to suppress immune responses against self-antigens or transplanted tissues
\item Tolerance-inducing drugs: Drugs designed to induce tolerance to specific antigens, rather than broadly suppressing the immune system
\end{itemize}

\section{Legacy}

The legacy of Frank Macfarlane Burnet and Peter Brian Medawar extends far beyond their Nobel Prize-winning work on acquired immunological tolerance. Their discoveries fundamentally changed our understanding of the immune system and laid the foundation for modern transplantation medicine.

\subsection{Foundation for Transplantation Medicine}

The work of Burnet and Medawar laid the foundation for modern transplantation medicine. The concept of acquired immunological tolerance has inspired the development of new strategies for preventing transplant rejection and has contributed to the success of organ transplantation as a life-saving procedure.

\subsection{Immunology as a Science}

Burnet and Medawar also contributed to the development of immunology as a distinct scientific discipline. Their work established the principles of modern immunology and inspired generations of researchers to study the immune system. The clonal selection theory remains one of a major theories in immunology and continues to guide research in the field.

\subsection{Inspiration for Future Researchers}

The work of Burnet and Medawar continues to inspire researchers in immunology and transplantation medicine. Their success in using elegant experiments to answer fundamental questions about the immune system demonstrates the power of basic research to improve human health. Their example has motivated generations of scientists to pursue careers in immunology and transplantation medicine.

\subsection{Recognition and Honors}

In addition to the Nobel Prize, Burnet and Medawar received numerous other honors and awards for their work. Burnet received the Order of Merit in 1958 and was elected to the Royal Society. Medawar received the Order of Merit in 1965 and was elected to the Royal Society.

Their contributions to immunology and medicine are remembered and celebrated by researchers and clinicians around the world. Their discoveries continue to influence our understanding of the immune system and have led to numerous advances in transplantation medicine and immunology, ensuring that their legacy will endure for generations to come.


\section{Modern Developments}

Burnet and Medawar's immunological tolerance work has evolved into modern transplant immunology and autoimmunity research. Understanding of T cell tolerance has led to tolerance-inducing therapies for organ transplantation (belatacept) and autoimmune diseases. Regulatory T cell (Treg) therapy is in clinical trials for type 1 diabetes and transplant rejection. Understanding of checkpoint inhibition (Nobel 2018) builds on their tolerance concepts.
